\section{Plano de Trabalho e Cronograma}

O estágio de pesquisa terá duração de 12 meses (Setembro de 2026 a Agosto de 2027) e será dividido nas seguintes fases principais:

\begin{enumerate}
    \item \textbf{Mês 1-2: Integração e Estudo (Fase 1)} \\
    Instalação na instituição anfitriã; familiarização com a infraestrutura \texttt{JuliaSmoothOptimizers} (JSO); estudo de métodos de região de confiança e operadores proximais estruturados.
    \item \textbf{Mês 3-5: Concepção Algorítmica Híbrida (Fase 2)} \\
    Desenvolvimento teórico de operadores proximais estruturados como pré-processadores globais e sua integração com os refinadores locais de descida de coordenadas e combinatórios estudados no Brasil. Análise preliminar de viabilidade.
    \item \textbf{Mês 6-8: Implementação Computacional e Relatório Parcial (Fase 3)} \\
    Implementação paralela dos algoritmos projetados na linguagem Julia, integrando os modelos desenvolvidos no Brasil ao framework JSO. Elaboração do Relatório Científico anual da FAPESP (fevereiro de 2027).
    \item \textbf{Mês 9-10: Benchmarking e Validação (Fase 4)} \\
    Execução de baterias de testes escaláveis em problemas de regressão linear e logística com penalização $L_0$ e esparsidade em grupos. Comparação contra pacotes estabelecidos na literatura.
    \item \textbf{Mês 11-12: Consolidação e Redação (Fase 5)} \\
    Consolidação dos resultados, integração das metodologias desenvolvidas na tese de doutorado e elaboração do Relatório Científico Final do BEPE exigido pela FAPESP.
\end{enumerate}

O fluxo de transferência do conhecimento se dará por meio de reuniões regulares de acompanhamento online com o orientador no Brasil, garantindo que todo o desenvolvimento de código e avanços teóricos sejam diretamente incorporados à tese de doutorado.
